% !TeX root = ../../ONMIPA Alin.tex
% Acuan: Zhang (2009), Bab 5, Soal 5.6--5.8, 5.19--5.23,
% 5.26--5.35, dan 5.50; Zhang (2011), bagian 1.4, khususnya
% Gram--Schmidt, komplemen ortogonal, dan ketaksamaan Bessel.
\chapter{Ortogonalitas dan Proyeksi}
\label{ch:ortogonalitas}

Basis ortonormal mengubah pencarian koordinat menjadi perhitungan
hasil kali dalam. Dari basis ini, kita memperoleh dekomposisi
ortogonal dan proyeksi, yang menjawab pertanyaan: vektor mana
dalam suatu subruang yang paling dekat dengan vektor yang diberikan?
Kecuali dinyatakan lain, ruang dalam bab ini berdimensi hingga.
Hasil kali dalam tetap linear pada argumen kedua.

\section{Sistem Ortogonal dan Ortonormal}

\begin{definisi}[][Daftar ortogonal dan ortonormal][ortho:def]
    Daftar $v_1,\ldots,v_m$ disebut ortogonal jika
    $\langle v_i,v_j\rangle=0$ untuk $i\ne j$.
    Daftar itu disebut ortonormal jika, sebagai tambahan,
    setiap vektornya bernorma satu; ekuivalen dengan
    $\langle v_i,v_j\rangle=\delta_{ij}$.
\end{definisi}

Setiap daftar ortogonal yang semua anggotanya tak nol bebas linear.
Memang, dari $\sum c_jv_j=0$, ambil hasil kali dalam dengan
$v_i$ pada argumen pertama untuk memperoleh $c_i\|v_i\|^2=0$.
Selain itu, ekspansi langsung memberikan
\[
    \left\|\sum_jc_jv_j\right\|^2=\sum_j|c_j|^2\|v_j\|^2.
\]
Syarat tak nol tidak boleh diabaikan dalam kesimpulan kebebasan linear.

\begin{teorema}[][Koordinat ortonormal dan identitas Parseval][ortho:parseval]
    Jika $e_1,\ldots,e_n$ basis ortonormal bagi $V$, maka
    \[
        x=\sum_{j=1}^n\langle e_j,x\rangle e_j,
        \qquad \|x\|^2=\sum_{j=1}^n|\langle e_j,x\rangle|^2.
    \]
    Lebih umum,
    \[
        \langle x,y\rangle
        =\sum_{j=1}^n\overline{\langle e_j,x\rangle}\langle e_j,y\rangle.
    \]
\end{teorema}
\begin{proof}
    Tulis $x=\sum c_je_j$. Mengambil $\langle e_i,x\rangle$
    memberi $c_i$. Ekspansi hasil kali dalam dua kombinasi
    tersebut hanya menyisakan suku dengan indeks yang sama,
    sehingga kedua identitas lainnya mengikuti.
\end{proof}

Jika $(v_1,\ldots,v_n)$ basis ortogonal, koefisien $v_j$ dalam
ekspansi $x$ adalah $\langle v_j,x\rangle/\|v_j\|^2$.
Penyebut boleh dihilangkan hanya ketika basisnya ortonormal.

\begin{teorema}[][Ketaksamaan Bessel][ortho:bessel]
    Jika $e_1,\ldots,e_m$ ortonormal, maka untuk setiap $x$,
    \[
        \sum_{j=1}^m|\langle e_j,x\rangle|^2\le\|x\|^2.
    \]
    Kesamaan berlaku tepat ketika
    $x\in\operatorname{span}\{e_1,\ldots,e_m\}$.
\end{teorema}
\begin{proof}
    Letakkan $p=\sum_j\langle e_j,x\rangle e_j$ dan $r=x-p$.
    Untuk setiap $i$, $\langle e_i,r\rangle=0$, sehingga $p\perp r$.
    Maka
    \[
        \|x\|^2=\|p\|^2+\|r\|^2
          =\sum_j|\langle e_j,x\rangle|^2+\|r\|^2.
    \]
    Kesamaan berlaku tepat ketika $r=0$.
\end{proof}

\section{Proses Gram--Schmidt}

\begin{teorema}[][Ortonormalisasi Gram--Schmidt][ortho:gram-schmidt]
    Untuk daftar bebas linear $v_1,\ldots,v_m$, definisikan
    \[
        u_1=v_1,\quad e_1=\frac{u_1}{\|u_1\|},\qquad
        u_k=v_k-\sum_{j=1}^{k-1}\langle e_j,v_k\rangle e_j,
        \quad e_k=\frac{u_k}{\|u_k\|}.
    \]
    Setiap $u_k$ tak nol, daftar $(e_1,\ldots,e_m)$ ortonormal,
    dan untuk setiap $k$ berlaku
    \[
        \operatorname{span}\{e_1,\ldots,e_k\}
        =\operatorname{span}\{v_1,\ldots,v_k\}.
    \]
\end{teorema}
\begin{proof}
    Gunakan induksi. Jika $e_1,\ldots,e_{k-1}$ sudah ortonormal,
    maka $\langle e_i,u_k\rangle=0$ untuk $i<k$.
    Jika $u_k=0$, vektor $v_k$ berada dalam ruang pembangun
    vektor-vektor sebelumnya, bertentangan dengan kebebasan linear.
    Normalisasi memberi vektor satuan baru yang ortogonal
    terhadap semua vektor sebelumnya.
    Rumus $u_k$ dan bentuk baliknya menyatakan bahwa kedua
    daftar membangun ruang yang sama pada setiap tahap.
\end{proof}

Akibatnya, setiap ruang hasil kali dalam berdimensi hingga mempunyai
basis ortonormal. Daftar ortonormal yang belum lengkap dapat
diperluas menjadi basis ortonormal: perluas dahulu menjadi basis
sembarang, lalu terapkan Gram--Schmidt. Vektor ortonormal yang
sudah ada pada awal daftar tidak berubah.

Jika daftar awal bergantung linear, suatu sisa $u_k$ dapat nol.
Sisa tersebut harus dilewati, bukan dibagi dengan normanya.
Proses pada vektor berikutnya tetap dapat dilanjutkan untuk
memperoleh basis ortonormal ruang pembangun daftar awal.

\begin{exam}[Ortonormalisasi tiga vektor real]
    Terapkan Gram--Schmidt pada
    $v_1=(1,1,0)^T$, $v_2=(1,0,1)^T$, $v_3=(0,1,1)^T$.
\end{exam}
\begin{solution}
    Langkah pertama memberi $e_1=(1,1,0)^T/\sqrt2$.
    Selanjutnya,
    \begin{align*}
        u_2&=v_2-\tfrac12(1,1,0)^T=(1/2,-1/2,1)^T,\\
        e_2&=(1,-1,2)^T/\sqrt6,\\
        u_3&=v_3-\tfrac12(1,1,0)^T-\tfrac16(1,-1,2)^T
             =\tfrac23(-1,1,1)^T.
    \end{align*}
    Jadi $e_3=(-1,1,1)^T/\sqrt3$.
    Ketiganya bernorma satu dan hasil kali setiap dua vektor
    berbeda nol, sehingga menjadi basis ortonormal $\R^3$.
\end{solution}

\begin{exam}[Basis ortonormal polinom]
    Pada $\mathcal P_{\le2}(\R)$ dengan
    $\langle p,q\rangle=\int_{-1}^1p(t)q(t)\,dt$,
    ortonormalkan daftar $(1,t,t^2)$.
\end{exam}
\begin{solution}
    Simetri interval membuat fungsi ganjil berintegral nol.
    Maka $1\perp t$, dengan norma kuadrat $2$ dan $2/3$.
    Sisa untuk $t^2$ adalah $t^2-1/3$, yang bernorma kuadrat $8/45$.
    Hasil akhirnya adalah
    \[
        e_0=\frac1{\sqrt2},\qquad
        e_1=\sqrt{\frac32}\,t,\qquad
        e_2=\sqrt{\frac58}\,(3t^2-1).
    \]
    Ortogonalitas di sini ditentukan oleh integral, bukan oleh
    hasil kali daftar koefisien dalam basis monomial.
\end{solution}

\section{Komplemen Ortogonal dan Dekomposisi}

\begin{definisi}[][Komplemen ortogonal][ortho:komplemen]
    Untuk $S\subseteq V$, definisikan
    \[
        S^\perp=\{x\in V:\langle s,x\rangle=0
            \text{ untuk setiap }s\in S\}.
    \]
\end{definisi}

Himpunan $S^\perp$ selalu subruang, bahkan jika $S$ bukan subruang.
Selain itu, $S^\perp=(\operatorname{span}S)^\perp$.
Jika $S\subseteq T$, maka $T^\perp\subseteq S^\perp$;
semakin banyak syarat ortogonal, semakin kecil ruang yang tersisa.

\begin{teorema}[][Dekomposisi ortogonal][ortho:dekomposisi]
    Untuk subruang $W$ dari ruang hasil kali dalam $V$ berdimensi $n$,
    \[
        V=W\oplus W^\perp,\qquad
        \dim W^\perp=n-\dim W,\qquad (W^\perp)^\perp=W.
    \]
    Untuk himpunan sembarang $S\subseteq V$,
    $(S^\perp)^\perp=\operatorname{span}S$.
\end{teorema}
\begin{proof}
    Pilih basis ortonormal $e_1,\ldots,e_r$ dari $W$ dan perluas
    menjadi basis ortonormal $e_1,\ldots,e_n$ dari $V$.
    Ekspansi koordinat menunjukkan bahwa $W^\perp$ dibangun
    oleh $e_{r+1},\ldots,e_n$. Ini membuktikan dekomposisi dan
    rumus dimensi. Vektor yang ortogonal terhadap semua anggota
    $W^\perp$ tepat kombinasi $e_1,\ldots,e_r$, sehingga
    komplemen dua kali mengembalikan $W$.
    Pernyataan terakhir mengikuti $S^\perp=(\operatorname{span}S)^\perp$.
\end{proof}

Untuk subruang $U,W$, berlaku
\[
    (U+W)^\perp=U^\perp\cap W^\perp,\qquad
    (U\cap W)^\perp=U^\perp+W^\perp.
\]
Identitas pertama mengikuti langsung dari definisi.
Untuk yang kedua, terapkan identitas pertama pada $U^\perp,W^\perp$
dan ambil komplemen ortogonal sekali lagi.
Penggunaan komplemen ganda di sini memakai asumsi dimensi hingga.

\begin{exam}[Ruang matriks bertrace nol]
    Dalam $\Mat_n(\C)$ dengan hasil kali dalam Frobenius,
    tentukan komplemen ortogonal dari
    $W=\{X:\tr X=0\}$ dan dekomposisi ortogonal setiap matriks $A$.
\end{exam}
\begin{solution}
    Karena $\langle I,X\rangle_F=\tr X$, diperoleh
    $W=\operatorname{span}\{I\}^\perp$ dan $W^\perp=\operatorname{span}\{I\}$.
    Dekomposisinya adalah
    \[
        A=\left(A-\frac{\tr A}{n}I\right)+\frac{\tr A}{n}I.
    \]
    Suku pertama bertrace nol dan suku kedua matriks skalar.
    Jarak $A$ ke $W$ adalah $|\tr A|/\sqrt n$.
\end{solution}

\section{Proyeksi dan Vektor Terdekat}

\begin{definisi}[][Proyeksi ortogonal][ortho:proyeksi]
    Jika $x=p+r$ dengan $p\in W$ dan $r\in W^\perp$,
    definisikan $P_Wx=p$.
    Operator $P_W$ disebut proyeksi ortogonal ke $W$.
\end{definisi}

Keunikan dekomposisi membuat definisi ini sah. Jika $(e_1,\ldots,e_r)$
basis ortonormal $W$, maka
\[
    P_Wx=\sum_{j=1}^r\langle e_j,x\rangle e_j.
\]
Rumus ini juga membuktikan linearitas, $P_W^2=P_W$,
$\operatorname{im}P_W=W$, dan $\ker P_W=W^\perp$.
Selain itu, $I-P_W=P_{W^\perp}$, serta
\[
    \|x\|^2=\|P_Wx\|^2+\|x-P_Wx\|^2.
\]
Dengan demikian $\|P_Wx\|\le\|x\|$, dengan kesamaan tepat ketika $x\in W$.

\begin{teorema}[][Karakterisasi pendekatan terbaik][ortho:terdekat]
    Untuk $p\in W$, pernyataan berikut ekuivalen:
    \begin{enumerate}[label=(\roman*),nosep]
        \item $p=P_Wx$;
        \item $x-p\in W^\perp$;
        \item $\|x-p\|\le\|x-w\|$ untuk semua $w\in W$.
    \end{enumerate}
    Vektor terdekat ini unik. Bahkan, untuk setiap $w\in W$,
    \[
        \|x-w\|^2=\|x-P_Wx\|^2+\|P_Wx-w\|^2.
    \]
\end{teorema}
\begin{proof}
    Dua pernyataan pertama adalah keunikan dekomposisi ortogonal.
    Untuk setiap $w\in W$, tulis $x-w=(x-P_Wx)+(P_Wx-w)$.
    Kedua suku ortogonal, sehingga Pythagoras memberi identitas
    yang dinyatakan. Nilainya minimum tepat ketika suku kedua
    nol, yaitu $w=P_Wx$. Ini membuktikan pernyataan ketiga dan keunikan.
\end{proof}

\begin{exam}[Proyeksi ke hiperbidang jumlah koordinat nol]
    Pada $\F^n$, letakkan $W=\{x:\sum_jx_j=0\}$.
    Tentukan $P_Wx$ dan jarak $x$ ke $W$.
\end{exam}
\begin{solution}
    Komplemen ortogonal $W$ dibangun oleh $\boldsymbol{1}=(1,\ldots,1)^T$.
    Dengan $\mu=\frac1n\sum_jx_j$, diperoleh
    \[
        P_Wx=x-\mu\boldsymbol{1},\qquad
        d(x,W)=\sqrt n\,|\mu|
          =\frac{|\sum_jx_j|}{\sqrt n}.
    \]
    Jadi proyeksi mengurangi setiap koordinat dengan rata-ratanya.
\end{solution}

\begin{exam}[Polinom linear terdekat dengan polinom kubik]
    Dalam $\mathcal P_{\le3}(\R)$ dengan hasil kali dalam integral
    pada $[-1,1]$, carilah polinom di $W=\operatorname{span}\{1,t\}$
    yang paling dekat dengan $t^3$.
\end{exam}
\begin{solution}
    Basis $(1,t)$ ortogonal. Koefisien proyeksi adalah
    \[
        \frac{\langle1,t^3\rangle}{\|1\|^2}=0,\qquad
        \frac{\langle t,t^3\rangle}{\|t\|^2}
          =\frac{2/5}{2/3}=\frac35.
    \]
    Jadi pendekatan terbaik adalah $3t/5$ dan galat kuadratnya
    $\|t^3\|^2-\|3t/5\|^2=2/7-6/25=8/175$.
\end{solution}

\section{Rumus Matriks dan Kuadrat Terkecil}

Jika $(w_1,\ldots,w_r)$ basis $W$ yang belum ortonormal,
tulis $P_Wx=\sum_jc_jw_j$. Syarat sisa ortogonal memberi
\[
    \sum_{j=1}^r\langle w_i,w_j\rangle c_j=\langle w_i,x\rangle,
    \qquad 1\le i\le r.
\]
Jadi $Gc=b$, dengan $G$ matriks Gram basis dan
$b_i=\langle w_i,x\rangle$. Karena basis bebas linear, $G$
positif definit dan $c=G^{-1}b$.

Untuk $W=\operatorname{im}A\subseteq\F^m$ dengan kolom-kolom
$A\in\Mat_{m\times r}(\F)$ bebas linear, rumus tersebut menjadi
\[
    P_W=A(A^*A)^{-1}A^*.
\]
Jika hasil kali dalamnya $x^*Hy$ dengan $H$ positif definit,
rumus yang sesuai adalah $P_W=A(A^*HA)^{-1}A^*H$.
Invers pada kedua rumus mensyaratkan kolom-kolom $A$ bebas linear.

\begin{exam}[Proyeksi berubah ketika hasil kali dalam berubah]
    Pada $\R^2$, gunakan $H=\begin{pmatrix}2&1\\1&2\end{pmatrix}$
    dan $\langle x,y\rangle_H=x^THy$.
    Tentukan proyeksi $(0,1)^T$ ke garis yang dibangun oleh $e_1$.
\end{exam}
\begin{solution}
    Koefisiennya adalah
    $\langle e_1,(0,1)^T\rangle_H/\langle e_1,e_1\rangle_H=1/2$.
    Maka proyeksinya $(1/2,0)^T$.
    Sisa $(-1/2,1)^T$ memang ortogonal terhadap $e_1$ dalam
    metrik $H$, meskipun tidak ortogonal dalam hasil kali dalam standar.
\end{solution}

\begin{teorema}[][Persamaan normal kuadrat terkecil][ortho:normal]
    Untuk $A\in\Mat_{m\times n}(\F)$ dan $b\in\F^m$, vektor
    $c$ meminimumkan $\|Ac-b\|$ jika dan hanya jika
    \[
        A^*Ac=A^*b.
    \]
    Solusi minimum selalu ada. Vektor $Ac$ yang dihasilkan unik,
    tetapi koefisien $c$ unik tepat ketika $\ker A=\{0\}$.
\end{teorema}
\begin{proof}
    Vektor $Ac$ terdekat dengan $b$ tepat ketika
    $b-Ac\perp\operatorname{im}A$.
    Ini ekuivalen dengan $A^*(b-Ac)=0$, yakni persamaan normal.
    Proyeksi ke $\operatorname{im}A$ selalu ada, dan menurut
    definisi citra proyeksi itu dapat ditulis sebagai $Ac$.
    Keunikan vektor proyeksi sudah diketahui. Jika $c_0$ salah
    satu solusi, seluruh solusi adalah $c_0+\ker A$.
\end{proof}

\begin{exam}[Garis yang paling sesuai dengan tiga data]
    Carilah garis $y=a+bt$ yang meminimumkan jumlah kuadrat
    galat pada data $(0,1)$, $(1,2)$, dan $(2,2)$.
\end{exam}
\begin{solution}
    Matriks desain dan vektor data adalah
    \[
        A=\begin{pmatrix}1&0\\1&1\\1&2\end{pmatrix},\qquad
        y=\begin{pmatrix}1\\2\\2\end{pmatrix}.
    \]
    Persamaan normal memberi
    \[
        \begin{pmatrix}3&3\\3&5\end{pmatrix}
        \begin{pmatrix}a\\b\end{pmatrix}
        =\begin{pmatrix}5\\6\end{pmatrix}.
    \]
    Maka $a=7/6$, $b=1/2$.
    Sisa $y-A(a,b)^T=(-1/6,1/3,-1/6)^T$ ortogonal terhadap
    kedua kolom $A$, dan jumlah kuadrat galat minimum adalah $1/6$.
\end{solution}

\section{Latihan Pilihan}

\begin{latihan}\label{ortho:lat-lengkapi}
    Lengkapi $u_1=(1,1,1,1)^T/2$ dan $u_2=(1,1,-1,-1)^T/2$
    menjadi basis ortonormal $\R^4$.
\end{latihan}

\begin{latihan}\label{ortho:lat-kompleks}
    Pada $\C^2$, letakkan $W=\operatorname{span}\{(1,i)^T\}$.
    Tentukan $W^\perp$, proyeksi $x=(1,1)^T$ ke $W$, dan jaraknya.
\end{latihan}

\begin{latihan}\label{ortho:lat-polinom}
    Lanjutkan ortonormalisasi $(1,t,t^2,t^3)$ terhadap hasil
    kali dalam integral pada $[-1,1]$. Berikan basis ortonormal
    lengkap bagi $\mathcal P_{\le3}(\R)$.
\end{latihan}

\begin{latihan}\label{ortho:lat-kuadratik}
    Tentukan real $a,b$ yang meminimumkan
    $\int_{-1}^1(t^2-a-bt)^2\,dt$, beserta nilai minimumnya.
\end{latihan}

\begin{latihan}\label{ortho:lat-kendala}
    Di antara $x\in\R^3$ yang memenuhi
    $x_1+x_2+x_3=1$ dan $x_1-x_3=2$, tentukan vektor bernorma
    terkecil dan norma kuadrat minimumnya.
\end{latihan}

\begin{latihan}\label{ortho:lat-dimensi}
    Dalam ruang berdimensi $n$, misalkan $u_1,\ldots,u_p$
    ortogonal dan semuanya tak nol. Jika $v_1,\ldots,v_q$
    masing-masing ortogonal terhadap semua $u_i$ dan $p+q>n$,
    buktikan bahwa daftar $v_1,\ldots,v_q$ bergantung linear.
\end{latihan}

\begin{latihan}\label{ortho:lat-dua-proyeksi}
    Misalkan $P,Q$ matriks proyeksi ortogonal pada $\C^n$.
    Buktikan bahwa $PQ$ merupakan matriks proyeksi ortogonal
    jika dan hanya jika $PQ=QP$. Dalam keadaan itu, tentukan citranya.
\end{latihan}

\begin{latihan}\label{ortho:lat-kontraksi}
    Misalkan $P^2=P$ pada ruang hasil kali dalam berdimensi hingga.
    Buktikan bahwa $P$ proyeksi ortogonal jika dan hanya jika
    $\|Px\|\le\|x\|$ untuk setiap $x$.
\end{latihan}

\ifpembahasan
\section{Pembahasan Latihan}

\subsection*{Latihan \ref{ortho:lat-lengkapi}: pola tanda yang saling meniadakan}
Ambil $u_3=(1,-1,1,-1)^T/2$ dan $u_4=(1,-1,-1,1)^T/2$.
Norma kuadrat masing-masing adalah empat kali $1/4$, yaitu satu.
Pada hasil kali dua vektor berbeda, terdapat dua suku $1/4$
dan dua suku $-1/4$, sehingga hasilnya nol.
Empat vektor ortonormal di $\R^4$ membentuk basis.

\subsection*{Latihan \ref{ortho:lat-kompleks}: konjugasi pada koefisien proyeksi}
Untuk $u=(1,i)^T$, persamaan $\langle u,(a,b)^T\rangle=a-ib=0$
memberi $W^\perp=\operatorname{span}\{(1,-i)^T\}$.
Karena $\langle u,x\rangle=1-i$ dan $\|u\|^2=2$,
\[
    P_Wx=\frac{1-i}{2}\begin{pmatrix}1\\i\end{pmatrix}
        =\begin{pmatrix}(1-i)/2\\(1+i)/2\end{pmatrix}.
\]
Norma kuadrat proyeksinya satu, sedangkan $\|x\|^2=2$.
Jadi jarak $x$ ke $W$ adalah satu.

\subsection*{Latihan \ref{ortho:lat-polinom}: pisahkan paritas polinom}
Tiga vektor pertama adalah $e_0,e_1,e_2$ pada contoh sebelumnya.
Karena $t^3$ ganjil, proyeksinya ke $e_0,e_2$ nol.
Sisa setelah mengurangi komponen pada $t$ adalah $t^3-3t/5$,
dengan norma kuadrat $8/175$. Jadi vektor terakhir dapat dipilih
\[
    e_3=\sqrt{\frac78}(5t^3-3t).
\]
Basis lengkapnya ialah
$1/\sqrt2$, $\sqrt{3/2}\,t$, $\sqrt{5/8}(3t^2-1)$,
dan $\sqrt{7/8}(5t^3-3t)$.

\subsection*{Latihan \ref{ortho:lat-kuadratik}: syarat sisa ortogonal}
Proyeksi $t^2$ ke $\operatorname{span}\{1,t\}$ mempunyai
koefisien $a=(2/3)/2=1/3$ dan $b=0$.
Nilai minimum adalah
\[
    \int_{-1}^1(t^2-1/3)^2\,dt=\frac25-\frac29=\frac8{45}.
\]
Keunikan proyeksi menjamin tidak ada pasangan lain yang mencapai minimum.

\subsection*{Latihan \ref{ortho:lat-kendala}: pilih solusi yang ortogonal terhadap kernel}
Dua vektor baris kendala, $(1,1,1)^T$ dan $(1,0,-1)^T$,
saling ortogonal. Solusi yang berada dalam ruang pembangun
keduanya adalah
\[
    x_0=\tfrac13(1,1,1)^T+(1,0,-1)^T
       =(4/3,1/3,-2/3)^T.
\]
Seluruh solusi lain berbentuk $x_0+t(1,-2,1)^T$.
Karena $x_0\perp(1,-2,1)^T$, norma kuadratnya adalah
$\|x_0\|^2+6t^2=7/3+6t^2$.
Minimum $7/3$ dicapai hanya pada $t=0$.

\subsection*{Latihan \ref{ortho:lat-dimensi}: hitung ruang yang masih tersedia}
Daftar $u_i$ bebas linear, sehingga ruang pembangunnya berdimensi $p$.
Semua $v_j$ berada dalam komplemen ortogonal yang berdimensi $n-p$.
Karena $q>n-p$, daftar tersebut pasti bergantung linear.

\subsection*{Latihan \ref{ortho:lat-dua-proyeksi}: simetri dan idempotensi}
Matriks proyeksi ortogonal memenuhi $P^*=P$ dan $P^2=P$,
dan hal yang sama berlaku untuk $Q$.
Sifat Hermitian dapat dilihat dari rumus
$P=EE^*$ dengan kolom-kolom $E$ basis ortonormal citranya.
Sebaliknya, matriks Hermitian idempoten mempunyai kernel
ortogonal terhadap citra, sehingga merupakan proyeksi ortogonal.
Jika $PQ$ proyeksi ortogonal, maka
$PQ=(PQ)^*=QP$.
Jika $PQ=QP$, maka $(PQ)^*=PQ$ dan $(PQ)^2=P^2Q^2=PQ$,
sehingga $PQ$ proyeksi ortogonal.
Citranya adalah $\operatorname{im}P\cap\operatorname{im}Q$:
setiap $PQx$ tetap oleh $P$ dan $Q$, dan setiap vektor pada
irisan tersebut tetap oleh $PQ$.

\subsection*{Latihan \ref{ortho:lat-kontraksi}: uji komponen kernel}
Arah perlu telah dibuktikan dengan Pythagoras.
Untuk arah cukup, ambil $u\in\operatorname{im}P$ dan $v\in\ker P$.
Untuk setiap skalar $t$, berlaku $P(u+tv)=u$, sehingga
$\|u\|^2\le\|u+tv\|^2$.
Jika $v\ne0$, pilih $t=-\langle v,u\rangle/\|v\|^2$.
Ekspansi norma kemudian memberi
\[
    \|u\|^2\le\|u\|^2-\frac{|\langle v,u\rangle|^2}{\|v\|^2},
\]
sehingga $\langle v,u\rangle=0$.
Kasus $v=0$ otomatis. Jadi citra dan kernel ortogonal.
Karena $P^2=P$, keduanya merupakan komplemen langsung;
maka $P$ tepat proyeksi ortogonal ke citranya.
\fi
